% 双臂电桥（开尔文电桥）原理电路图
% 拓扑：R1(Px1→B)、R2(B→Ps2) 与主回路 Rx-r-Rs 并联；
%       R3(Px2→D)、R4(D→Ps1) 与粗导线 r 并联；检流计 G 接在 B、D 之间。
%       平衡时由 R1 I1 = Rx I3 + R3 I2、R2 I1 = Rs I3 + R4 I2 得 Rx = (R1/R2) Rs。
\newcommand{\ArmBridge}{%
\begin{circuitikz}[european,line width=0.7pt]

  % ---------- 顶部主回路：四端电阻 Rx —— 粗导线 r —— 比较臂 Rs ----------
  \draw (1.5,4.8) rectangle (5.1,5.2);
  \node[font=\small] at (3.3,5.0) {$R_x$};
  \draw (6.4,4.8) rectangle (10.0,5.2);
  \node[font=\small] at (8.2,5.0) {$R_s$};
  \draw[line width=2pt] (5.1,5.0) -- (6.4,5.0);
  \node[below,font=\footnotesize] at (5.75,4.94) {$r$};

  % ---------- 比率臂：R1/R2、R3/R4 各留出间隙 ----------
  \draw (2.1,4.8) -- (2.1,4.4);
  \draw (2.1,4.4) to[R] (2.1,3.4);
  \draw (2.1,3.4) -- (2.1,2.8);
  \draw (2.1,2.8) to[R] (2.1,1.8);
  \draw (4.5,4.8) -- (4.5,4.4);
  \draw (4.5,4.4) to[R] (4.5,3.4);
  \draw (4.5,3.4) -- (4.5,2.8);
  \draw (4.5,2.8) to[R] (4.5,1.8);

  % ---------- 检流计 ----------
  \node[circle,draw,inner sep=1pt,font=\small] (G) at (3.3,3.1) {$G$};
  \draw (2.1,3.1) -- (G.west);
  \draw (G.east) -- (4.5,3.1);
  \fill (2.1,3.1) circle (1pt);
  \fill (4.5,3.1) circle (1pt);

  % ---------- 电位端回路 ----------
  \draw (2.1,1.8) -- (2.1,0.6) -- (9.4,0.6) -- (9.4,4.8);   % R2 → Ps2
  \draw (4.5,1.8) -- (4.5,1.2) -- (7.0,1.2) -- (7.0,4.8);   % R4 → Ps1

  % ---------- 直流电源 ----------
  \draw (1.5,5.0) -- (-0.4,5.0) -- (-0.4,-1.3) -- (4.0,-1.3);
  \draw (4.0,-1.52) -- (4.0,-1.08);
  \draw (4.45,-1.42) -- (4.45,-1.18);
  \draw (4.45,-1.3) -- (10.8,-1.3) -- (10.8,5.0) -- (10.0,5.0);
  \node[below,font=\small] at (4.22,-1.72) {$\varepsilon$};

  % ---------- 端子与元件标注 ----------
  \node[above left,font=\footnotesize]  at (1.45,5.06)  {$C_{x1}$};
  \node[above,font=\footnotesize]       at (4.95,5.22)  {$C_{x2}$};
  \node[above,font=\footnotesize]       at (6.55,5.22)  {$C_{s1}$};
  \node[above right,font=\footnotesize] at (10.05,5.06) {$C_{s2}$};
  \node[right,font=\footnotesize]       at (2.16,4.6)   {$P_{x1}$};
  \node[right,font=\footnotesize]       at (4.56,4.6)   {$P_{x2}$};
  \node[right,font=\footnotesize]       at (7.06,4.6)   {$P_{s1}$};
  \node[right,font=\footnotesize]       at (9.46,4.6)   {$P_{s2}$};
  \node[left,font=\footnotesize]        at (1.86,3.9)   {$R_1$};
  \node[left,font=\footnotesize]        at (1.86,2.3)   {$R_2$};
  \node[right,font=\footnotesize]       at (4.74,3.9)   {$R_3$};
  \node[right,font=\footnotesize]       at (4.74,2.3)   {$R_4$};
  \node[left,font=\footnotesize]        at (2.0,3.1)    {$B$};
  \node[right,font=\footnotesize]       at (4.6,3.1)    {$D$};

\end{circuitikz}}
